\section*{第 \ref{ch:g3:levers-and-scales} 章——杠杆与天平}

\begin{solution}{exo:g3:levers-and-scales:1}
跷跷板：木板是\omterm{def:g3:levers-and-scales:lever}{杠杆}，中间的支座是\omterm{def:g3:levers-and-scales:lever}{支点}。
撬棍：铁棒是\omterm{def:g3:levers-and-scales:lever}{杠杆}，圆木是\omterm{def:g3:levers-and-scales:lever}{支点}。剪刀：每一片带把手的刀片
都是\omterm{def:g3:levers-and-scales:lever}{杠杆}，两者的\omterm{def:g3:levers-and-scales:lever}{支点}都在中间的螺丝处。
\end{solution}

\begin{solution}{exo:g3:levers-and-scales:2}
离\omterm{def:g3:levers-and-scales:lever}{支点}越远：用力的位置越靠外，乘积就越大
——远处的小力，胜过近处的大石头。
\end{solution}

\begin{solution}{exo:g3:levers-and-scales:3}
$3 \times 4 = 12$；\ $2 \times 5 = 10$；\ $1 \times 6 = 6$。这三个乘积
互不相等，因此任意两个都不能互相平衡。
\end{solution}

\begin{solution}{exo:g3:levers-and-scales:4}
左边的乘积是 $1 \times 8 = 8$。两枚硬币需要 $8 = 2 \times
4$：标 $4$。四枚硬币需要 $8 = 4 \times 2$：标 $2$。
\end{solution}

\begin{solution}{exo:g3:levers-and-scales:5}
左边：$3 \times 3 = 9$。右边：$2 \times 4 = 8$。左边的
乘积更大，因此左边下降。
\end{solution}

\begin{solution}{exo:g3:levers-and-scales:6}
两臂等长时，托盘水平就说明两边重量相等，这样的比较是公平的。
若左臂更长，左边的重物就会乘以更大的
距离，于是每次称量，\omterm{def:g3:levers-and-scales:scale}{天平}都
会偏向左边。
\end{solution}

\begin{solution}{exo:g3:levers-and-scales:7}
袋子的重量恰好等于 $6$ 块积木的重量。拿走一块积木后，
装袋子的托盘变重，\omterm{def:g3:levers-and-scales:scale}{天平}向
袋子那一侧倾斜。
\end{solution}

\begin{solution}{exo:g3:levers-and-scales:8}
汤姆的乘积：$2 \times 3 = 6$。他的妹妹需要 $1 \times 6 = 6$：
她必须坐在 $6$ 号刻度处。
\end{solution}

\begin{solution}{exo:g3:levers-and-scales:9}
\omterm{def:g3:levers-and-scales:lever}{支点}是锤头压在木头上的位置。你的手
在锤柄的远端用力——距离长，因此乘积大——
而钉子的阻力离\omterm{def:g3:levers-and-scales:lever}{支点}很近：
\omterm{def:g3:levers-and-scales:lever}{杠杆}把你轻松的推力变成了巨人般的拉力。
\end{solution}

\begin{solution}{exo:g3:levers-and-scales:10}
平衡需要 $4 \times (\text{爷爷的刻度}) = 1 \times
(\text{佐伊的刻度})$。两种方法：爷爷在 $1$ 处、佐伊在 $4$ 处
（$4 = 4$），或爷爷在 $2$ 处、佐伊在 $8$ 处（$8 = 8$）。
\end{solution}

\begin{solution}{exo:g3:levers-and-scales:11}
更轻。在长臂一端，即使很轻的重物也会产生很大的乘积，因此
只要货物比短臂上诚实的积木\emph{轻}，托盘就会保持水平。
空秤检查能揭穿他：托盘里什么也不放时，
他的\omterm{def:g3:levers-and-scales:scale}{天平}就已经向长臂一侧倾斜。
\end{solution}
